%! TeX program = lualatex
% A numbered, empty grid for the curve-sketching exercise that closes the
% derivative tests lesson.  The window is cut for that exercise's data: the
% vertical asymptote at x = 1, the horizontal asymptote y = 2, the local
% maximum at (-2,4) and the inflection point near x = -4 all have to sit
% inside it with room left over, and the branch running down to -infinity has
% to have somewhere to go.
%
% The ticks keep their numbers, unlike the blank grids elsewhere in these
% notes: the data names particular points, so a student has to be able to
% find them.
\documentclass[tikz, crop]{standalone}
\input{../typography.tex.preamble}
\input{../tikz.tex.preamble}

\begin{document}

% ---------------------------------------------------------------------
% Page 1 -- the exercise built from a list of facts about f, f' and f''.
\begin{tikzpicture}
  \begin{axis}[
    basic curve,
    width=5in,
    xmin=-6.5, xmax=6.5,
    ymin=-4.5, ymax=6.5,
    ]
  \end{axis}
\end{tikzpicture}

% ---------------------------------------------------------------------
% Page 2 -- the final exam question, f(x) = x e^{-x^2/2}.  Everything the
% earlier parts ask students to find:
%
%   x-intercept         0, the only one
%   local minimum       x = -1,          f = -1/sqrt(e) = -0.607
%   local maximum       x =  1,          f =  1/sqrt(e) =  0.607
%   inflection points   x = 0 and +/- sqrt3 = +/- 1.732, where f = +/- 0.387
%   horizontal asymptote  y = 0, both directions
%
% The curve is odd, so the window is symmetric about the origin.  x runs to
% +/- 4 because that is where the Gaussian factor has flattened the curve
% onto its asymptote -- f(4) = 0.0013 -- and a student has to be able to see
% that happen rather than guess it.
%
% Equal units are kept, which is what fixes the proportions: the curve is
% eight units wide and never leaves the band between -0.61 and 0.61, so it
% genuinely is this flat, and a grid that stretched it vertically would
% teach the wrong shape.  What the flatness costs is precision -- 0.607 and
% 0.387 both fall inside one unit square -- so the grid is halved with minor
% lines, giving a 0.5 rule to read those heights against.
\begin{tikzpicture}
  \begin{axis}[
    basic curve,
    width=6.2in,
    xmin=-4.2, xmax=4.2,
    ymin=-2.1, ymax=2.1,
    minor tick num=1,
    grid=both,
    minor grid style={gray!12},
    ]
  \end{axis}
\end{tikzpicture}
\end{document}
